\section{Recovery of angular loss by a common physical collar}
\label{sec:v71-pooled-collar}
A chart may lose angular occupation while another chart of the same
exact record gains it. The original comparison treats the losses
separately. We retain its controlled source and use the remaining
receiving capacity before declaring a loss. The operation is a scalar
weighting of the original probability space. It neither joins orbits
nor exchanges their integer labels.

\subsection{A positive recovery principle}
\begin{lemma}[Maximal recovery with a prescribed capacity]
\label{lem:v71-capacity-recovery}
Let $X,O,Y,S,D$ be nonnegative finite numbers, with
\[
 O\le\min(X,Y),\qquad S\le E O,\qquad D\le E(X-O),
 \qquad E>0.
\]
For $K\ge1$, define
\begin{equation}\label{eq:v71-recovery-coefficient}
 \theta_K=
 \begin{cases}
 \min\{1,(KY-O)/(X-O)\},&X>O,\\
 1,&X=O.
 \end{cases}
\end{equation}
Then $0\le\theta_K\le1$ and
\begin{align}
 S_K:=S+\theta_KD&\le E\min(X,KY),
                                      \label{eq:v71-capacity-controlled}\\
 D_K:=(1-\theta_K)D&\le E(X-KY)_+,
 \qquad S_K+D_K=S+D.                 \label{eq:v71-capacity-deficit}
\end{align}
Moreover $S_K\ge S$ and $0\le D_K\le D$. When $X>O$,
$\theta_K$ is the largest $\theta\in[0,1]$ satisfying
$O+\theta(X-O)\le KY$. For fixed data, $D_K$ is nonincreasing
in $K$.
\end{lemma}
\begin{proof}
Since $KY\ge O$, the coefficient is nonnegative. Directly,
\[
 O+\theta_K(X-O)=\min(X,KY),\qquad
 (1-\theta_K)(X-O)=(X-KY)_+.
\]
These identities include $X=O$, when $D=0$ and $KY\ge X$.
The bounds for $S,D$ prove the two estimates. The remaining statements
follow from the defining inequality and from monotonicity in $K$.
\end{proof}

The optimality here is for the displayed scalar capacity constraint.
It is not an optimal-transport statement about the physical measure.
In particular unequal guards or physical Jacobians have not been
replaced by equal ones; $E$ bounds them in the source comparison.

\subsection{Exactly labelled capacities at a common roof}
Fix $m,R$, $\lambda=(n,k)$, $0<\chi\le1/10$ and
$0<\varepsilon\le\chi/4$. Set
\[
 h=h_\chi,\qquad j=h/2,\qquad J=(0,j),\qquad
 E_\chi=\exp(K_\chi r_\chi).
\]
Every angular profile below has the same exact label $\lambda$.
For a selected chart put
\[
 \overline L_{z,\lambda}=j^{-1}\int_0^jL_{z,\lambda}(v)\,dv,
 \qquad
 o_{z,\lambda}(u)=j^{-1}\int_0^j
                  \min\{L_{z,\lambda}(u),L_{z,\lambda}(v)\}\,dv.
\]
Here $L_{z,\lambda}(u)=\ell_{z,\lambda}(\sqrt{2u})$ is the
actual angular fraction, including every physical and section test.
For a roof $t$ write
\[
 \mathcal Z_\lambda(t)=\{z:0<u_z(t):=t-t_z<h\}.
\]
The selected chart set depends on $m,R$; it is finite at each count.
Define the three capacities
\begin{equation}\label{eq:v71-pooled-capacities}
 \begin{split}
 \mathsf X_\lambda(t)&=\sum_{z\in\mathcal Z_\lambda(t)}
                         J_zL_{z,\lambda}(u_z(t)),\\
 \mathsf O_\lambda(t)&=\sum_{z\in\mathcal Z_\lambda(t)}
                         J_zo_{z,\lambda}(u_z(t)),\\
 \mathsf Y_\lambda(t)&=\sum_{z\in\mathcal Z_\lambda(t)}
                         J_z\overline L_{z,\lambda}.
 \end{split}
\end{equation}
The dependence on $m,R,\chi,\varepsilon$ is suppressed, not averaged. The
normalization $c_*^{-1}$ is already included in each $J_z$. There
is no summation over different return indices or displacements.

Use the unchanged $s_{70},d_{70}$ of
\eqref{eq:v70-complete-source-identity}. On common coarea versions,
\begin{equation}\label{eq:v71-physical-capacity-bounds}
 0\le\mathsf O_\lambda\le\min(\mathsf X_\lambda,\mathsf Y_\lambda),
 \quad s_{70}\le E_\chi\mathsf O_\lambda,
 \quad d_{70}\le E_\chi(\mathsf X_\lambda-\mathsf O_\lambda).
\end{equation}
Indeed $b_z(t_z+u)\le E_\chi J_zL_{z,\lambda}(u)$, since the
original guard is at most one. Multiply by the old overlap and loss
weights, and use \eqref{eq:v70-capacity-identity}. This also proves
that zero angular capacity produces zero source; no quotient is taken
there. These estimates hold for the actual nonconstant guards.

Insert $(\mathsf X_\lambda,\mathsf O_\lambda,\mathsf Y_\lambda)$
in \eqref{eq:v71-recovery-coefficient}, and write the resulting
measurable coefficient as $\theta_{K,m,R,\lambda}(t)$. Define
\begin{equation}\label{eq:v71-physical-recovery}
 s_{71,K}=s_{70}+\theta_{K,m,R,\lambda}d_{70},\qquad
 d_{71,K}=(1-\theta_{K,m,R,\lambda})d_{70}.
\end{equation}
This is a weighting on the original source: on an actual chart state
$x$ use $\theta_{K,m,R,\lambda}(L_{m,R}(x))$ in its original
loss weight. Every multiplier lies in $[0,1]$. The densities in
\eqref{eq:v71-physical-recovery} therefore are the pushforwards of
these source weights, not merely an algebraic replacement of densities.
The weights may depend measurably on the whole exactly labelled
profile; no spectral norm of that insertion is used.

\begin{theorem}[Ordered recovery on the complete selected disks]
\label{thm:v71-pooled-height}
Fix $K\ge1$, $\chi$ and $\varepsilon$ before the collision limit.
The complete original clearance source has the positive identity
\begin{equation}\label{eq:v71-complete-source}
 b^{\varepsilon,\mathrm{clr}}=s_{71,K}+d_{71,K}+e_{70}.
\end{equation}
The outside source is exactly the one in
\eqref{eq:v66-complete-positive-partition}. On every selected disk,
including its outer annuli,
\begin{align}
 s_{71,K}&\ge s_{70},\qquad 0\le d_{71,K}\le d_{70},
                                      \label{eq:v71-monotone-recovery}\\
 d_{71,K}(t)&\le E_\chi
              (\mathsf X_\lambda(t)-K\mathsf Y_\lambda(t))_+.
                                      \label{eq:v71-pooled-deficit}
\end{align}
Moreover
\begin{equation}\label{eq:v71-pooled-height}
 \mathcal H(s_{71,K})
 \le \frac32 CK h_\chi
          \exp\{(1+1/\sqrt2)K_\chi r_\chi\}
 \le C'K\chi^6.
\end{equation}
The constants are independent of the collision count, the exact label
and the number of charts. Selected-chart and outside pieces are
disjoint events; controlled and residual chart pieces are
complementary weights, not disjoint events.
\end{theorem}
\begin{proof}
Apply Lemma~\ref{lem:v71-capacity-recovery} to
\eqref{eq:v71-physical-capacity-bounds}. It gives the pointwise
claims and $s_{71,K}\le E_\chi K\mathsf Y_\lambda$.
Adding the unchanged outside source proves the complete identity.
For the ordered estimate, the lower unguarded physical profile on
$0<v<j$ gives
\[
 J_z\overline L_{z,\lambda}
       \le e^{K_\chi\sqrt{2j}}\widehat q_{z,j;\lambda}.
\]
At roof $t$ all contributing centers belong to $[t-h,t]$.
Thus the sum of these positive collar averages is bounded by
Theorem~\ref{thm:v69-comparison-collar}, with $w=h$ and $s=j$.
The interval length plus collar width is $h+j=3h/2$; the receiving
average has already been divided by $j$. This proves the first
bound, since $\sqrt{2j}=r_\chi/\sqrt2$. The exponential is
uniformly bounded because $K_\chi r_\chi=Cr_0\chi$ and
$h_\chi=r_0^2\chi^6/2$.

The estimate is applied to one positive physical event after the
sum over centers. In particular the dependence of $\theta$ on $t$
or on all the charts introduces no chart-count or label-count factor.
All widths remain fixed throughout that collision limit.
\end{proof}

\begin{corollary}[Complete selected-source height with a receiving reserve]
\label{cor:v71-reserve-class}
Suppose $\mathsf X_\lambda(t)\le K\mathsf Y_\lambda(t)$ for
almost every roof, uniformly in $R,\lambda$ at all sufficiently
large counts. Then the entire original selected-chart source has
height at most $C'K\chi^6$. No individual angular profile is required
to be monotone, or to have a bounded variation independent of the
word. In particular the old loss $d_{70}$ has this same height bound.
\end{corollary}
\begin{proof}
The residual in \eqref{eq:v71-pooled-deficit} vanishes. The entire
selected source equals $s_{71,K}$, and $d_{70}$ is its positive
subsource. Apply \eqref{eq:v71-pooled-height}.
\end{proof}

This is a condition on the actual, physically weighted capacities;
it is not asserted for every Lorentz record. It is strictly weaker
than requiring each chart to have no loss to its receiving collar.
For example, as a scalar diagnostic, take two equal weights and equal
centers with profiles $L_1(u)=\one_{(j,h)}(u)$ and
$L_2(u)=\one_{(0,j)}(u)$. At $u>j$ the first profile is a late
birth and its old source is wholly loss, whereas
$\mathsf X=\mathsf Y$ and the new residual is zero. More generally
complementary profiles with equal weights have this property,
regardless of their number of radial changes. These are capacity
examples, not assertions that arbitrary profiles occur as Lorentz
words. Different labels, radii or collision counts cannot be pooled.
For an isolated late birth with $\mathsf Y=0<\mathsf X$, no finite
$K$ recovers its source. No exhaustion claim is hidden in the parameter
$K$.
